\documentclass[11pt,a4paper,bibliography=totoc]{scrartcl}
\usepackage[utf8]{inputenc}
\usepackage[german]{babel}
\usepackage[T1]{fontenc}
\usepackage{amsmath}
\usepackage{amssymb}
\usepackage{mathtools}
\usepackage{booktabs}
\usepackage{microtype}
\usepackage{geometry}
\usepackage{hyperref}

\geometry{a4paper, left=25mm, right=25mm, top=30mm, bottom=30mm}

\title{\textsf{Arithmetische Vakuumpolarisation und EABC-Phasenräume}}
\subtitle{Ein Forschungsprogramm zu Primzahlvierlingen, Tschebyscheff-Bias, modularer Dämpfung und diskreter Symmetriebrechung \\ \vspace{0.5em} \large Antragsskizze für ein dreijähriges Forschungsprojekt mit fünf wissenschaftlichen Mitarbeitern}
\author{\textbf{Thomas Hofbauer} \\ Bamberg}
\date{\today}

\begin{document}

\maketitle

\begin{abstract}
\noindent Dieses Forschungsprogramm schlägt die systematische Untersuchung eines arithmetischen Phasenmodells vor, das Primzahlen nicht nur als isolierte Elemente der natürlichen Zahlen, sondern als Bestandteile diskreter Symmetrie-, Schalen- und Übergangsstrukturen betrachtet. Im Zentrum steht das sogenannte EABC-Modell: die vier primzahlfähigen Restklassen modulo 12,
\[
E \equiv 1, \qquad A \equiv 5, \qquad B \equiv 7, \qquad C \equiv 11 \pmod{12},
\]
sowie die besondere Rolle von Primzahlvierlingen
\[
Q(p) = (p, p+2, p+6, p+8).
\]
Für $p > 3$ durchläuft ein solcher Vierling im EABC-System genau zwei chirale Ordnungen: $ABCE$ oder $CEAB$. Diese beiden Orientierungen bilden den Ausgangspunkt für eine präzise Untersuchung arithmetischer Phasen-Bias-Effekte. Das Programm verbindet klassische Fragen der analytischen Zahlentheorie, insbesondere Tschebyscheff-Bias und Primzahlrennen, mit geometrischen Schalenmodellen, modularen Dämpfungsparametern wie $q = e^{-\pi}$, und spektralen Operatoransätzen im Sinne einer diskreten, arithmetischen Symmetriebrechung.

Ziel des Projektes ist nicht die vorschnelle Behauptung einer physikalischen Identität zwischen Zahlentheorie und Quantenfeldtheorie, sondern die Entwicklung eines mathematisch kontrollierten Modells, in dem Begriffe wie „Vakuumpolarisation“, „Sombrero-Potential“ und „Phasenbias“ als präzise definierte arithmetische Strukturen auftreten. Das beantragte dreijährige Projekt mit fünf wissenschaftlichen Mitarbeitern soll klären, ob das EABC-Modell einen eigenständigen Beitrag zu den Bereichen Primzahlrennen, admissible prime constellations, modulare Dämpfungsstrukturen, diskrete Schalenmodelle und spektrale Operatoren leisten kann.
\end{abstract}

\tableofcontents
\vfill
\newpage

\section{Historischer Hintergrund}

\subsection{Von Euklid zum Fundamentalsatz der Arithmetik}
Die Geschichte der Primzahlen beginnt mit Euklids Beweis der Unendlichkeit der Primzahlen. Der spätere Fundamentalsatz der Arithmetik bildet eine der Grundsäulen der Zahlentheorie:
\[
n = \prod_{i=1}^{r} p_i^{\alpha_i},
\]
wobei die Darstellung bis auf die Reihenfolge eindeutig ist. Primzahlen sind damit die irreduziblen Bausteine der natürlichen Zahlen. Jede strukturelle Aussage über ihre Verteilung berührt das Fundament der Arithmetik selbst.

Das hier vorgeschlagene Modell respektiert diesen klassischen Ausgangspunkt. Es behauptet nicht, Primzahlvierlinge seien die Bausteine aller Zahlen. Vielmehr werden Primzahlvierlinge als besonders symmetrische lokale Konfigurationen innerhalb des durch den Fundamentalsatz garantierten Primzahlgitters betrachtet.

\subsection{Dirichlet, Riemann und die Verteilung in Restklassen}
Dirichlets Satz über Primzahlen in arithmetischen Progressionen zeigt, dass für teilerfremde Restklassen $a$ modulo $q$ unendlich viele Primzahlen $p \equiv a \pmod q$ existieren. Der Primzahlsatz für arithmetische Progressionen verfeinert diese Aussage asymptotisch:
\[
\pi(x; q, a) \sim \frac{1}{\varphi(q)} \operatorname{Li}(x).
\]
Für modulo 12 bedeutet dies, dass die vier primzahlfähigen Klassen $1, 5, 7, 11 \pmod{12}$ asymptotisch gleich dicht sind. Doch die Gleichverteilung im Hauptterm schließt sekundäre Asymmetrien nicht aus. Genau hier treten der Tschebyscheff-Bias und Primzahlrennen auf.

\subsection{Tschebyscheff-Bias und Primzahlrennen}
Der klassische Tschebyscheff-Bias beschreibt die Beobachtung, dass Primzahlen der Form $4n+3$ über weite Bereiche häufiger „vorne liegen“ als Primzahlen der Form $4n+1$. Formal betrachtet man die Differenz
\[
\Delta_4(x) = \pi(x; 4, 3) - \pi(x; 4, 1).
\]
Asymptotisch verschwindet der relative Hauptunterschied, dennoch ist die Verteilung der Vorzeichen von $\Delta_4(x)$ nicht neutral. Diese Erscheinung wurde im 20.\ Jahrhundert durch Arbeiten über Primzahlrennen, explizite Formeln und Nullstellen von Dirichlet-$L$-Funktionen präzisiert.

Eine entscheidende arithmetische Quelle dieses Bias liegt darin, dass Quadrate ungerader Zahlen modulo 4 stets in der Klasse 1 liegen: $p^2 \equiv 1 \pmod 4$. Daraus ergeben sich sekundäre Verzerrungen, die in der expliziten Formel sichtbar werden. Das vorliegende Projekt verallgemeinert diese Perspektive auf ein vierphasiges modulo-12-System.

\subsection{Primzahlvierlinge als minimale symmetrische Primkonstellationen}
Primzahlzwillinge, Primzahldrillinge und Primzahlvierlinge gehören zu den klassischen \emph{admissible prime constellations}. Ein Primzahlvierling hat die Form
\[
Q(p) = (p, p+2, p+6, p+8).
\]
Die inneren Abstände sind $(2, 4, 2)$. Sein Mittelpunkt ist $m = p+4$, sodass sich die Konstellation symmetrisch schreiben lässt als:
\[
Q(p) = (m-4, m-2, m+2, m+4).
\]
Diese Darstellung zeigt, dass ein Primzahlvierling eine diskrete Spiegelstruktur besitzt. Er ist nicht nur eine zufällige Ansammlung von vier Primzahlen, sondern eine lokale symmetrische Zelle innerhalb der natürlichen Zahlen.

Für $p > 3$ gilt im EABC-System:
\begin{align*}
p \equiv 5 \pmod{12} \quad &\Rightarrow \quad Q(p) \sim ABCE, \\
p \equiv 11 \pmod{12} \quad &\Rightarrow \quad Q(p) \sim CEAB.
\end{align*}
Damit besitzt jeder Primzahlvierling eine chirale Orientierung.

\section{Grundidee des EABC-Modells}

\subsection{Die vier primzahlfähigen Klassen modulo 12}
Alle Primzahlen größer als 3 liegen in genau einer der vier Restklassen $1, 5, 7, 11 \pmod{12}$. Wir bezeichnen diese Klassen durch
\[
E \equiv 1, \qquad A \equiv 5, \qquad B \equiv 7, \qquad C \equiv 11 \pmod{12}.
\]
Damit entsteht ein natürlicher vierphasiger arithmetischer Raum $\mathcal{H}_{\mathrm{EABC}} \cong \mathbb{C}^4$ mit den Basisvektoren $e_E, e_A, e_B, e_C$. Ein arithmetischer Zustand kann formal geschrieben werden als
\[
z = z_E e_E + z_A e_A + z_B e_B + z_C e_C.
\]
Die Norm
\[
|z|^2 = |z_E|^2 + |z_A|^2 + |z_B|^2 + |z_C|^2
\]
beschreibt eine abstrakte Phasenenergie.

\subsection{Chirale Vierlingszellen}
Für einen Primzahlvierling $Q(p) = (p, p+2, p+6, p+8)$ ergeben sich modulo 12 genau zwei zulässige Orientierungen: $ABCE$ oder $CEAB$. Daraus definieren wir die Zählfunktionen
\begin{align*}
N_{ABCE}(x) &= \# \{p \leq x : Q(p) \text{ ist Primzahlvierling der Orientierung } ABCE\}, \\
N_{CEAB}(x) &= \# \{p \leq x : Q(p) \text{ ist Primzahlvierling der Orientierung } CEAB\}.
\end{align*}
Der zugehörige chirale Bias ist
\[
\Delta_{\mathrm{EABC}}(x) = N_{ABCE}(x) - N_{CEAB}(x).
\]
Die normierte Form lautet
\[
\mathcal{B}_{\mathrm{EABC}}(x) = \frac{N_{ABCE}(x) - N_{CEAB}(x)}{N_{ABCE}(x) + N_{CEAB}(x)}.
\]
Diese Größe ist einer der zentralen Untersuchungsgegenstände des Projekts.

\subsection{Arithmetische Vakuumpolarisation}
Wir verwenden den Begriff „arithmetische Vakuumpolarisation“ in einem präzise eingeschränkten Sinn. Er bezeichnet nicht eine physikalische Vakuumpolarisation. Vielmehr beschreibt er die Tatsache, dass ein formal symmetrisches arithmetisches System eine statistisch asymmetrische Phasenstruktur besitzen kann. Das zugrundeliegende Schema lautet:
\[
\text{asymptotische Gleichverteilung} + \text{sekundäre Korrekturterme} \quad \Rightarrow \quad \text{beobachtbarer Bias}.
\]
Beim klassischen Tschebyscheff-Bias betrifft dies Restklassen modulo 4. Im EABC-Modell betrifft es chirale Übergangsstrukturen modulo 12, insbesondere die beiden Vierlingsorientierungen $ABCE$ und $CEAB$.

\section{Das Schalenmodell}

\subsection{Quadratschalen}
Wir zerlegen die natürlichen Zahlen in Quadratschalen:
\[
\mathcal{S}_n = \{m \in \mathbb{N} : n^2 \leq m < (n+1)^2\}.
\]
Die Länge der $n$-ten Schale ist
\[
|\mathcal{S}_n| = (n+1)^2 - n^2 = 2n + 1.
\]
Diese Schalen wachsen linear. Dadurch entsteht eine natürliche nichtlineare Skalenstruktur. Die Quadrate bilden die Schalenränder: $n^2$ und $(n+1)^2$.

\subsection{Pronische Mittellinien}
Zwischen zwei Quadraten liegt die pronische Zahl $P_n = n(n+1)$. Sie erfüllt $n^2 < P_n < (n+1)^2$. Explizit gilt $P_n = n^2 + n$. Die inverse pronische Skala ist
\[
n(P) = \frac{-1 + \sqrt{1 + 4P}}{2}.
\]
Damit kann eine große Zahl $P$ auf ihre effektive Schalenkoordinate zurückgeführt werden. Im Modell übernehmen Quadrate und pronische Zahlen verschiedene Rollen:
\begin{align*}
\text{Quadrate} &= \text{Schalenränder}, \\
\text{pronische Zahlen} &= \text{Schalenachsen}.
\end{align*}

\subsection{Primzahlvierlinge innerhalb von Schalen}
Für jede Schale definieren wir
\[
\mathcal{Q}_n = \{p : Q(p) \subseteq \mathcal{S}_n\}.
\]
Das heißt, es gilt $p \geq n^2$ und $p+8 < (n+1)^2$. Die zentrale Frage lautet: Wie verteilen sich die Orientierungen $ABCE$ und $CEAB$ innerhalb der Quadratschalen? Dazu definieren wir lokale Schalen-Biasgrößen:
\[
\mathcal{B}_n = \frac{N_{ABCE}(\mathcal{S}_n) - N_{CEAB}(\mathcal{S}_n)}{N_{ABCE}(\mathcal{S}_n) + N_{CEAB}(\mathcal{S}_n)},
\]
sofern der Nenner nicht verschwindet. Diese Größen erlauben eine feinere Analyse als globale Zählfunktionen allein.

\section{Sombrero-Potential und Symmetriebrechung}

\subsection{Higgs-artiges Potential als mathematische Analogie}
Das Higgs-Potential der Physik besitzt die Form $V(\phi) = \lambda \left(|\phi|^2 - v^2\right)^2$. Seine Minima liegen nicht bei $\phi = 0$, sondern auf der Menge $|\phi| = v$. Wir übertragen diese Struktur nicht physikalisch, sondern formal auf den EABC-Phasenraum. Für $z \in \mathbb{C}^4$ definieren wir ein symmetrisches EABC-Potential:
\[
V_0(z) = \lambda \left(|z|^2 - v^2\right)^2.
\]
Dieses Potential besitzt eine Schale von Minima bei $|z| = v$.

\subsection{Bias-Term}
Um die arithmetische Symmetriebrechung zu modellieren, ergänzen wir einen Bias-Term:
\[
V_{\mathrm{EABC}}(z) = \lambda \left(|z|^2 - v^2\right)^2 + \varepsilon B_{\mathrm{EABC}}(z).
\]
Hier ist $B_{\mathrm{EABC}}(z)$ kein physikalisches Feld, sondern eine mathematische Funktion, die aus Restklassenstatistiken, Vierlingsorientierungen oder Dirichlet-Charakteren konstruiert werden kann. Ein diskreter Kandidat ist der globale Bias $\mathcal{B}_{\mathrm{EABC}}(x)$. Ein glatterer Kandidat kann durch gewichtete Summen definiert werden:
\[
B_{\mathrm{EABC}}(x; w) = \frac{\sum_{p \leq x} w(p)\chi_{\mathrm{chir}}(p)}{\sum_{p \leq x} w(p)},
\]
wobei
\[
\chi_{\mathrm{chir}}(p) = 
\begin{cases}
+1, & \text{falls } Q(p) \text{ Orientierung } ABCE \text{ hat}, \\
-1, & \text{falls } Q(p) \text{ Orientierung } CEAB \text{ hat}, \\
0,  & \text{falls } Q(p) \text{ kein Primzahlvierling ist}.
\end{cases}
\]

\subsection{Potentialkuhlen für Vierlingsphasen}
Die beiden Orientierungen $ABCE$ und $CEAB$ können als zwei diskrete Potentialkuhlen interpretiert werden. Ein Primzahlvierling entspricht dann einer arithmetischen Phasenauskristallisierung in einer dieser beiden Kuhlen. Das Modell behauptet nicht, dass Primzahlen physikalisch in einem Potential rollen. Es behauptet nur: Die statistischen Bias-Effekte können formal wie Symmetriebrechungen in einem diskreten Potentialsystem beschrieben werden. Diese Unterscheidung ist wesentlich für die mathematische Seriosität des Programms.

\section{Modulare Dämpfung und der Parameter \texorpdfstring{$e^{-\pi}$}{e\textasciicircum(-\textbackslash pi)}}

\subsection{Motivation}
Der Parameter $q = e^{-\pi}$ tritt in vielen modularen und elliptischen Kontexten natürlich auf. Numerisch gilt $e^{-\pi} \approx 0{,}0432139$. Im vorliegenden Modell wird $q$ nicht als Naturkonstante vorausgesetzt, sondern als kanonische kleine modulare Skala untersucht. Wir betrachten mögliche Entwicklungen der Form
\[
F(q) = c_0 + c_1 q + c_2 q^2 + c_3 q^3 + \cdots.
\]
Für $q = e^{-\pi}$ sind die höheren Terme stark gedämpft.

\subsection{Anwendung auf EABC-Bias}
Ein möglicher Forschungsansatz lautet:
\[
\mathcal{B}_{\mathrm{EABC}}(x) \approx c_0(x) + c_1(x)e^{-\pi} + c_2(x)e^{-2\pi} + \cdots.
\]
Dabei ist offen, ob eine solche Entwicklung empirisch stabil oder theoretisch begründbar ist. Das Projekt wird daher drei Fragen untersuchen:
\begin{enumerate}
    \item Gibt es in den EABC-Biasdaten stabile modulare Skalen?
    \item Tritt $e^{-\pi}$ tatsächlich als natürlicher Dämpfungsfaktor auf?
    \item Lässt sich ein solcher Faktor aus Thetafunktionen, Dirichlet-$L$-Funktionen oder elliptischen Strukturen herleiten?
\end{enumerate}

\section{Hurwitz-Strukturen, Quaternionen und Oktonionen}

\subsection{Quaternionische Vierfachstruktur}
Die vier EABC-Klassen legen eine vierdimensionale Struktur nahe. Eine natürliche algebraische Vergleichsstruktur sind die Quaternionen $\mathbb{H}$. Hurwitz-Quaternionen bilden eine besonders symmetrische ganzzahlige Struktur innerhalb der Quaternionenalgebra. Sie besitzen enge Verbindungen zu Gittern, Normformen und diskreten Symmetriegruppen. Eine vorsichtige Forschungsfrage lautet daher: Kann das EABC-System als arithmetische Vierphasenstruktur mit quaternionischen Symmetrien modelliert werden?

\subsection{Oktonionische Erweiterung}
Oktonionen $\mathbb{O}$ sind achtdimensional und nichtassoziativ. Sie können formal als Verdopplung quaternionischer Strukturen verstanden werden: $\mathbb{O} \cong \mathbb{H} \oplus \mathbb{H}$. Für das EABC-Modell bietet sich eine Interpretation an als
\[
(E, A, B, C) \quad + \quad (E', A', B', C').
\]
Die zweite Vierergruppe könnte duale, gespiegelte oder torsional verschobene Phasen herleiten. Die Nichtassoziativität der Oktonionen wäre dann nicht bloß ein technisches Hindernis, sondern möglicherweise ein Modell für arithmetische Torsion: $(ab)c - a(bc) \neq 0$. Die zentrale Frage lautet: Gibt es messbare arithmetische Effekte in Vierlingsübergängen, die sich als Nichtassoziativitäts- oder Torsionsphänomene modellieren lassen?

\section{Spektrale Operatoren}

\subsection{Motivation aus Hilbert-Pólya}
Die Hilbert-Pólya-Idee fragt, ob die nichttrivialen Nullstellen der Riemannschen Zetafunktion als Spektrum eines selbstadjungierten Operators verstanden werden können. Das vorliegende Projekt erhebt nicht den Anspruch, eine Lösung der Riemannschen Vermutung zu liefern. Es nimmt jedoch die spektrale Perspektive ernst und fragt: Welche Operatoren entstehen natürlich aus EABC-Übergängen, Primzahlvierlingen und Schalenstrukturen?

\subsection{EABC-Übergangsoperator}
Wir definieren einen Übergangsoperator $T : \mathbb{C}^4 \to \mathbb{C}^4$ durch zyklische Verschiebung $T(E,A,B,C) = (A,B,C,E)$. In Matrixform:
\[
T = \begin{pmatrix}
0 & 0 & 0 & 1 \\
1 & 0 & 0 & 0 \\
0 & 1 & 0 & 0 \\
0 & 0 & 1 & 0
\end{pmatrix}.
\]
Die Eigenwerte sind $1, i, -1, -i$. Damit besitzt der EABC-Raum eine natürliche vierphasige Spektralstruktur.

\subsection{Chiraler Operator}
Die beiden Orientierungen $ABCE$ und $CEAB$ legen einen chiralen Operator nahe:
\[
\Gamma_{\chi} e_{ABCE} = +e_{ABCE}, \qquad \Gamma_{\chi} e_{CEAB} = -e_{CEAB}.
\]
Der Bias kann dann als Erwartungswert geschrieben werden:
\[
\mathcal{B}_{\mathrm{EABC}}(x) = \langle \Gamma_{\chi} \rangle_x.
\]
Dies macht die Analogie zu quantenmechanischen Erwartungswerten formal präzise, ohne eine physikalische Gleichsetzung zu behaupten.

\subsection{Schalen-Hamiltonian}
Für Quadratschalen definieren wir einen diskreten Hamiltonian:
\[
H_n = H_{\mathrm{radial}}(n) + H_{\mathrm{EABC}}(n) + H_{\mathrm{bias}}(n).
\]
Dabei steht $H_{\mathrm{radial}}(n)$ für die Schalenhöhe, $H_{\mathrm{EABC}}(n)$ für Übergänge zwischen EABC-Klassen und $H_{\mathrm{bias}}(n)$ für chirale Asymmetrien. Ein einfacher ansatz lautet:
\[
H_{\mathrm{radial}}(n) = \sqrt{n} \cdot I, \qquad H_{\mathrm{bias}}(n) = \varepsilon \mathcal{B}_n \Gamma_{\chi}.
\]
Ziel ist nicht ein willkürliches Modell, sondern die systematische Suche nach Operatoren, deren Spektren reproduzierbare Korrelationen mit bekannten arithmetischen Daten zeigen.

\section{Forschungsfragen}
Das Projekt soll folgende Hauptfragen beantworten:
\begin{description}
    \item[Frage 1: Gibt es stabile EABC-Bias-Effekte?] Untersucht werden die Größen $\mathcal{B}_{\mathrm{EABC}}(x)$ und $\mathcal{B}_n$ für große Werte von $x$ und $n$. Zu klären ist, ob beobachtete Asymmetrien statistisch signifikant, stabil und von bekannten Primzahlrennen erklärbar sind.
    \item[Frage 2: Welche Rolle spielen Quadratschalen und pronische Achsen?] Es soll geprüft werden, ob Primzahlvierlinge innerhalb der Schalen $[n^2, (n+1)^2]$ nichttriviale Verteilungsmuster zeigen. Besonders relevant sind die Abstände zu $n^2$, $n(n+1)$ und $(n+1)^2$.
    \item[Frage 3: Gibt es eine modulare Dämpfungsskala?] Empirisch und theoretisch soll untersucht werden, ob Parameter wie $e^{-\pi}$ in Bias-, Spektral- oder Übergangsstatistiken natürlich auftreten.
    \item[Frage 4: Gibt es eine quaternionische oder oktonionische Modellierung?] Zu prüfen ist, ob die EABC-Struktur sinnvoll mit Hurwitz-Quaternionen, Gitterstrukturen oder nichtassoziativen oktonionischen Erweiterungen verbunden werden kann.
    \item[Frage 5: Lassen sich natürliche Operatoren konstruieren?] Das Projekt soll eine Familie von EABC-Operatoren entwickeln und deren Spektren untersuchen. Insbesondere sollen Vergleiche mit Referenzspektren erfolgen (Zufallsmatrizen, Zeta- und $L$-Nullstellen, Graphenspektren).
\end{description}

\section{Methodik}

\subsection{Analytische Zahlentheorie}
Verwendet werden klassische Methoden: Primzahlsatz in arithmetischen Progressionen, Dirichlet-Charaktere, explizite Formeln, Primzahlrennen sowie Hardy-Littlewood-Heuristiken für Primkonstellationen. Insbesondere soll geprüft werden, ob der EABC-Bias vollständig aus bekannten expliziten Formeln ableitbar ist.

\subsection{Numerische Hochleistungsanalyse}
Es werden effiziente Algorithmen zur Erzeugung und Analyse von Primzahlvierlingen bis zu sehr großen Schranken entwickelt. Berechnet werden $N_{ABCE}(x)$, $N_{CEAB}(x)$, $\mathcal{B}_{\mathrm{EABC}}(x)$ und $\mathcal{B}_n$, sowie Verteilungen relativ zu Quadratschalen und pronischen Achsen, Gap-Statistiken und Spektralstatistiken.

\subsection{Spektralanalyse}
Für die Operatoren $H_n$, $T$ und $\Gamma_\chi$ werden Eigenwertverteilungen, Spacing-Statistiken und Korrelationen untersucht. Der zentrale Test lautet, ob die beobachteten Spektren robust gegenüber Modellvariationen sind.

\subsection{Algebraische Modellierung}
Die EABC-Struktur soll systematisch mit den algebraischen Objekten $\mathbb{Z}/12\mathbb{Z}$, $\mathbb{C}^4$, $\mathbb{H}$, den Hurwitz-Quaternionen sowie $\mathbb{O}$ verglichen werden.

\section{Arbeitsprogramm für drei Jahre}

\noindent\begin{tabular}{p{0.25\textwidth} p{0.7\textwidth}}
\toprule
\textbf{Zeitraum} & \textbf{Arbeitspakete und Meilensteine} \\
\midrule
\textbf{Jahr 1:} & \textbf{AP 1:} Mathematische Formalisierung (EABC-Phasenraum, Funktionale). \\
\emph{Fundament \& Daten} & \textbf{AP 2:} Numerische Datenbank (Python/C++/Julia-Pipeline, Schalenzuordnung). \\
                 & \textbf{AP 3:} Vergleich mit klassischen Primzahlrennen (modulo 4 vs.\ modulo 12). \\
                 & \emph{Ergebnisse:} Offener Datensatz, Preprints 1 \& 2. \\
\midrule
\textbf{Jahr 2:} & \textbf{AP 4:} Schalenanalyse (Verteilung um pronische Achsen, Signifikanztests). \\
\emph{Operatoren \& Dämpfung} & \textbf{AP 5:} EABC-Operatoren (Schalen-Hamiltonians, Spektralanalyse Software). \\
                 & \textbf{AP 6:} Modulare Dämpfung (Fit-Tests gegen $e^{-\pi}$, Thetafunktionen). \\
                 & \emph{Ergebnisse:} Operatoren-Bibliothek, Preprints 3, 4 \& 5. \\
\midrule
\textbf{Jahr 3:} & \textbf{AP 7:} Algebraische Erweiterungen (Hurwitz-Strukturen, Oktonionen-Torsion). \\
\emph{Synthese}  & \textbf{AP 8:} Spektralvergleich (Zeta-Nullstellen, Random-Matrix-Baselines). \\
                 & \textbf{AP 9:} Synthese und finale Abschlussmonographie. \\
                 & \emph{Ergebnisse:} Preprints 6 \& 7, Monographie. \\
\bottomrule
\end{tabular}

\section{Personalstruktur}
Beantragt wird eine Arbeitsgruppe von fünf wissenschaftlichen Mitarbeitern (jeweils 100\,\% TV-L E13) über drei Jahre:
\begin{itemize}
    \item \textbf{Stelle 1: Analytische Zahlentheorie} (Primzahlrennen, L-Funktionen, Heuristiken).
    \item \textbf{Stelle 2: Numerische Zahlentheorie} (HPC, Algorithmen, Datenbankaufbau, Pipelines).
    \item \textbf{Stelle 3: Spektraltheorie und Operatoren} (Eigenwertstatistiken, Zufallsmatrizen).
    \item \textbf{Stelle 4: Algebra und Geometrie} (Quaternionen, Hurwitz-Gitter, Oktonionen).
    \item \textbf{Stelle 5: Mathematische Physik und Modellkritik} (Potentialmodelle, Validierung).
\end{itemize}
Die fünfte Stelle übernimmt explizit die Rolle einer internen Kontrollinstanz, um eine unkontrollierte Physikalisierung rein mathematischer Analogien zu verhindern.

\section{Risikoanalyse \& Wissenschaftlicher Mehrwert}
Das Projekt besitzt den Charakter grundlagenorientierter Risikoforschung. Es ist strukturiert falsifizierbar, was seinen wissenschaftlichen Wert sichert:
\begin{itemize}
    \item \textbf{Risiko 1 (Klassische Erklärbarkeit):} Sollte der EABC-Bias vollständig klassisch über bekannte Dirichlet-Charaktere erklärt werden, liefert das Projekt eine wertvolle Einordnung eines neuen Phasenraums in die bestehende Theorie.
    \item \textbf{Risiko 2 (Instabilität von $e^{-\pi}$):} Ein negatives Resultat bezüglich der modularen Dämpfungshypothese schützt das Feld vor numerologischen Fehlinterpretationen.
    \item \textbf{Risiko 3 (Fehlende Operatorenstruktur):} Ohne robuste Spektralstrukturen verbleibt das Modell als reines kombinatorisch-statistisches Phänomen, was den Suchraum für zukünftige Operatormodelle präzise eingrenzt.
\end{itemize}

\section{Zentrale Hypothese}

\begin{center}
\renewcommand{\arraystretch}{1.5}
\begin{tabular}{|p{0.9\textwidth}|}
\hline
\textbf{Zentrale Arbeitshypothese des Forschungsprogramms:} \\
Die vier primzahlfähigen Restklassen modulo 12 bilden einen natürlichen EABC-Phasenraum. Primzahlvierlinge erscheinen darin als chirale Zellen mit den Orientierungen $ABCE$ und $CEAB$. Die statistische Asymmetrie dieser Orientierungen kann als arithmetischer Phasen-Bias untersucht werden. Quadratschalen, pronische Achsen, modulare Dämpfung und spektrale Operatoren liefern zusätzliche mathematische Strukturen, deren Tragfähigkeit systematisch und falsifizierbar geprüft werden soll. \\
\hline
\end{tabular}
\end{center}

\section{Schlussbemerkung}
Das EABC-Modell führt in seinem gegenwärtigen Zustand mehrere klassische Stränge der Mathematik zusammen:
\[
\text{Restklassen} + \text{Primzahlvierlinge} + \text{Bias} + \text{Schalen} + \text{Spektren} + \text{Symmetriebrechung}.
\]
Das beantragte dreijährige Forschungsprogramm präsentiert keine fertige Theorie, sondern eine ernsthafte, historisch begründete und methodisch kontrollierte Forschungsagenda zur Etablierung einer präzisen arithmetischen Phasengeometrie.

\end{document}